\chapter*{INTRODUCCIÓN}
\addcontentsline{toc}{chapter}{INTRODUCCIÓN}

El acceso oportuno a la información clínica es un pilar fundamental de la atención médica moderna. Un sistema de red local bien diseñado permite que los historiales clínicos electrónicos, imágenes radiológicas DICOM y registros de farmacia fluyan de forma inmediata entre los profesionales de salud, reduciendo tiempos de diagnóstico y minimizando errores derivados del uso de documentación en papel.

La presente memoria técnica documenta de manera exhaustiva y sistematizada el proyecto de diseño de la infraestructura física de red LAN del \textbf{Policlínico de Especialidades Médicas ``Salud \& Vida'' S.A.C.}, una edificación de \textbf{450~m$^2$ de planta única} ubicada en la ciudad de Tacna. El proyecto ha sido desarrollado como trabajo integrador de la asignatura \textbf{Redes I (Unidad I: Infraestructura de Red)}, cubriendo desde la definición de requerimientos organizacionales hasta el protocolo de certificación y la confección manual de la muestra física de cableado.

\section*{Objetivo General}

Diseñar la infraestructura de red de área local (LAN) del Policlínico ``Salud \& Vida'' S.A.C., bajo los estándares internacionales ANSI/TIA, garantizando conectividad Gigabit a todos los puestos de trabajo y disponibilidad de red inalámbrica Wi-Fi~6 en toda la planta, dentro de un presupuesto optimizado y con los márgenes técnicos requeridos para una operación de misión crítica médica.

\section*{Objetivos Específicos}

\begin{enumerate}[leftmargin=*, label=\arabic*.]
    \item Identificar y clasificar todos los ambientes, dispositivos y requerimientos de conectividad de la organización.
    \item Seleccionar y justificar la topología de red, los modelos de referencia y los medios de transmisión más apropiados para el entorno hospitalario.
    \item Diseñar el sistema de cableado estructurado bajo los estándares ANSI/TIA-568, 569, 606 y 607, incluyendo la distribución del gabinete rack 42U y el mapa de parcheo de los 56 puertos proyectados.
    \item Elaborar los planos físicos de distribución de tomas, canalizaciones y cobertura inalámbrica, junto con el cuadro de materiales valorizado (BOM).
    \item Definir el protocolo de construcción de cables UTP y el plan de pruebas de certificación, preparando al equipo para la defensa técnica individual ante el docente.
\end{enumerate}

\section*{Alcance y Limitaciones}

El alcance del proyecto comprende únicamente la \textbf{Unidad I: Infraestructura Física y Cableado Estructurado}. El direccionamiento IP, enrutamiento, seguridad lógica y servicios de red (DHCP, DNS, VPN) corresponden a unidades posteriores de la asignatura y quedan explícitamente fuera del alcance del presente documento.

\section*{Estructura del Documento}

La memoria se organiza en cinco capítulos, que corresponden a las cinco etapas del proyecto:

\begin{itemize}[leftmargin=*]
    \item \textbf{Capítulo 1:} Definición de la organización, inventario de ambientes y matriz de requerimientos de red.
    \item \textbf{Capítulo 2:} Selección de topología, aplicación de los modelos OSI/TCP-IP y justificación de los medios de transmisión.
    \item \textbf{Capítulo 3:} Ingeniería de cableado estructurado, diseño del rack 42U y cómputo de materiales.
    \item \textbf{Capítulo 4:} Planos de distribución física, diseño de canalizaciones, cobertura Wi-Fi~6 y BOM valorizado.
    \item \textbf{Capítulo 5:} Construcción de cable UTP bajo norma T568B y plan de pruebas de certificación.
\end{itemize}
